%!
% @file 04_implementation.tex
% @brief Implementation chapter for WordSprint report.

\chapter{Triển khai}
\label{chap:implementation}

% -------------------------------------------------
% Môi trường phát triển và quy trình build
% -------------------------------------------------
\section{Môi trường phát triển và quy trình build}
\label{sec:build_environment}
Dự án WordSprint được phát triển trên hệ điều hành Windows x64 với các công cụ sau:
\begin{itemize}
    \item \textbf{Compiler}: MSVC (đi kèm Visual Studio 2022) với chuẩn C++20.
    \item \textbf{CMake}: phiên bản ≥ 3.22 để cấu hình dự án và quản lý phụ thuộc Qt.
    \item \textbf{Qt}: Qt 6.11.3 (MSVC 2022 64‑bit) – các module Core, Gui, Qml, Quick, WebEngineQuick, WebChannel.
    \item \textbf{Node.js}: phiên bản 22.x và \textbf{pnpm} để quản lý phụ thuộc frontend.
    \item \textbf{Vite}: công cụ xây dựng nhanh cho React 19.
    \item \textbf{Tailwind CSS 4}: cho việc tạo kiểu UI.
\end{itemize}

\subsection{Cấu hình CMake (tập tin \texttt{CMakeLists.txt})}
\label{subsec:cmake}
\begin{lstlisting}[caption={CMakeLists.txt tại gốc dự án}, label={lst:cmake_root}]
cmake_minimum_required(VERSION 3.22)

project(WordSprint LANGUAGES CXX)

set(CMAKE_CXX_STANDARD 20)
set(CMAKE_CXX_STANDARD_REQUIRED ON)

if(NOT CMAKE_PREFIX_PATH AND EXISTS "C:/Qt/6.11.3/msvc2022_64")
    list(APPEND CMAKE_PREFIX_PATH "C:/Qt/6.11.3/msvc2022_64")
endif()

add_subdirectory(backend)
\end{lstlisting}
Trong file này, CMake tự động thêm đường dẫn Qt khi tồn tại, sau đó đưa thư mục \texttt{backend} vào quá trình biên dịch.

\subsection{Xây dựng backend}
\label{subsec:backend_build}
\begin{enumerate}
    \item Mở "Developer PowerShell for VS 2022".
    \item Kiểm tra công cụ: \texttt{cl /Bv}, \texttt{cmake --version}.
    \item Tạo thư mục build và chạy CMake:
    \begin{lstlisting}[caption={Câu lệnh build backend}, label={lst:backend_build}]
mkdir build
cd build
cmake -S .. -B . -G "Visual Studio 17 2022" -A x64 -DCMAKE_PREFIX_PATH="C:/Qt/6.11.3/msvc2022_64"
cmake --build . --config Release
\end{lstlisting}
    \item Sau khi biên dịch, tệp thực thi \texttt{WordSprint.exe} nằm trong \texttt{build/backend/Release}.
\end{enumerate}

\subsection{Xây dựng frontend}
\label{subsec:frontend_build}
\begin{enumerate}
    \item Mở một cửa sổ PowerShell trong thư mục \texttt{frontend}.
    \item Cài đặt các phụ thuộc:
    \begin{lstlisting}[caption={Cài đặt phụ thuộc frontend}, label={lst:frontend_deps}]
pm install -g pnpm   # nếu chưa có
pnpm install --frozen-lockfile
\end{lstlisting}
    \item Chạy máy chủ phát triển (dev):
    \begin{lstlisting}[caption={Khởi động Vite dev server}, label={lst:frontend_dev}]
pnpm dev
\end{lstlisting}
    \item Để tạo bản build tĩnh, dùng:
    \begin{lstlisting}[caption={Build frontend}, label={lst:frontend_build}]
pnpm build
\end{lstlisting}
    \item Sau khi build, các file tĩnh được đặt trong \texttt{frontend/dist} và được Qt WebEngine tải khi chạy chế độ Release.
\end{enumerate}

\subsection{Khởi chạy hệ thống}
\label{subsec:run}
\begin{enumerate}
    \item Chạy backend trong chế độ debug (tự động kết nối tới Vite dev server):
    \begin{lstlisting}[caption={Chạy ứng dụng trong debug}, label={lst:run_debug}]
build\backend\Release\WordSprint.exe --debug
\end{lstlisting}
    \item Hoặc chạy bản release, sẽ tải tài nguyên tĩnh từ \texttt{frontend/dist}:
    \begin{lstlisting}[caption={Chạy bản release}, label={lst:run_release}]
build\backend\Release\WordSprint.exe
\end{lstlisting}
\end{enumerate}

% -------------------------------------------------
% Thực hiện các chức năng cốt lõi
% -------------------------------------------------
\section{Thực hiện các chức năng cốt lõi}
\label{sec:core_features}
Dưới đây là các đoạn mã tiêu biểu mô tả cách triển khai các chức năng chính.

\subsection{Khởi tạo Qt và tải giao diện web}
\label{subsec:qt_init}
\begin{lstlisting}[caption={main.cpp – khởi tạo Qt và WebEngineView}, label={lst:main_cpp}]
int main(int argc, char *argv[]) {
    // Đặt flag Chromium để bật GPU
    qputenv("QTWEBENGINE_CHROMIUM_FLAGS",
            "--ignore-gpu-blocklist --enable-gpu-rasterization --enable-zero-copy");

    QCoreApplication::setAttribute(Qt::AA_ShareOpenGLContexts);
    QtWebEngineQuick::initialize();

    QGuiApplication app(argc, argv);
    app.setApplicationName("WordSprint");
    app.setOrganizationName("WordSprint");

    // Phân tích các tùy chọn dòng lệnh (debug / url)
    QCommandLineParser parser;
    parser.addHelpOption();
    parser.addVersionOption();
    parser.addOption({{"d", "debug"}, "Enable debug mode to live preview from Vite dev server."});
    parser.addOption({{"u", "url"}, "Custom URL to load in WebEngineView." , "url"});
    parser.process(app);

    bool isDebug = parser.isSet("debug");
    QUrl initialUrl;
    if (parser.isSet("url")) {
        initialUrl = QUrl::fromUserInput(parser.value("url"));
    } else if (isDebug) {
        const QString devUrl = qEnvironmentVariable("VITE_DEV_SERVER_URL", "http://localhost:5173");
        initialUrl = QUrl(devUrl);
    } else {
        // Tải file tĩnh đã build
        const QString distPath = QCoreApplication::applicationDirPath() + "/../frontend/dist/index.html";
        initialUrl = QUrl::fromLocalFile(QFileInfo(distPath).canonicalFilePath());
    }

    QQmlApplicationEngine engine;
    engine.setInitialProperties({{"initialUrl", initialUrl}});
    engine.rootContext()->setContextProperty("initialUrl", initialUrl);
    engine.load(QUrl("qrc:/qt/qml/WordSprint/qml/Main.qml"));
    return app.exec();
}
\end{lstlisting}
Đoạn mã trên quyết định URL nguồn (dev server hoặc tĩnh) và khởi tạo \texttt{QQmlApplicationEngine} để nhúng giao diện React.

\subsection{Main.qml – cửa sổ Qt WebEngineView}
\label{subsec:main_qml}
\begin{lstlisting}[caption={Main.qml – WebEngineView}, label={lst:main_qml}]
import QtQuick
import QtQuick.Window
import QtWebEngine

Window {
    id: root
    visible: true
    title: qsTr("WordSprint")
    property url initialUrl: "http://localhost:5173"

    WebEngineView {
        anchors.fill: parent
        url: root.initialUrl
        settings.accelerated2dCanvasEnabled: true
        settings.webGLEnabled: true
        settings.localContentCanAccessRemoteUrls: true
        settings.javascriptCanAccessClipboard: true
    }
}
\end{lstlisting}
WebEngineView chịu trách nhiệm hiển thị toàn bộ SPA React.

\subsection{Frontend – chuyển đổi hành động người dùng thành yêu cầu C++}
\label{subsec:frontend_bridge}
Trong React, khi người dùng khởi động trò chơi, một thông điệp JSON được gửi qua WebChannel. Đoạn mã mẫu minh họa cách gọi hàm C++:
\begin{lstlisting}[caption={Gửi lệnh startGame tới backend}, label={lst:frontend_js}]
// frontend/src/App.jsx (đoạn thu gọn)
function startGame() {
    const payload = { action: "startGame" };
    window.qt?.webChannel?.objects?.backend?.processMessage(JSON.stringify(payload));
}
\end{lstlisting}
Backend nhận thông điệp và thực hiện \texttt{GameEngine::startGame()} (được đăng ký trong C++ – không có trong ví dụ hiện tại nhưng logic tương tự). Các phản hồi được trả về dưới dạng JSON và cập nhật UI thông qua state React.

% -------------------------------------------------
% Kiểm thử chức năng và biên
% -------------------------------------------------
\section{Kiểm thử chức năng và biên}
\label{sec:testing}
Kiểm thử được thực hiện theo hai cấp: unit test cho backend (có thể dùng \texttt{Qt Test}) và end‑to‑end test cho frontend (công cụ \texttt{Playwright} hoặc \texttt{Cypress}).

\subsection{Kịch bản kiểm thử chức năng}
\label{subsec:test_cases}
\begin{itemize}
    \item \textbf{TC01 – Khởi động ứng dụng}: Kiểm tra \texttt{WordSprint.exe} chạy, WebEngineView tải URL đúng (dev hoặc tĩnh).
    \item \textbf{TC02 – Bắt đầu trò chơi}: Người dùng nhấn nút "Start", frontend gửi \texttt{startGame}. Backend trả về puzzle đầu tiên và thời gian giới hạn.
    \item \textbf{TC03 – Kiểm tra đáp án đúng}: Nhập đáp án đúng, backend trả về \texttt{correct=true} và cập nhật điểm theo \eqref{eq:score}.
    \item \textbf{TC04 – Kiểm tra đáp án sai}: Nhập đáp án sai, backend trả về \texttt{correct=false}, điểm không thay đổi, thời gian giảm.
    \item \textbf{TC05 – Hết thời gian}: Khi đồng hồ đếm ngược tới 0, backend gửi \texttt{gameOver} và lưu hồ sơ.
    \item \textbf{TC06 – Lưu và tải hồ sơ}: Kiểm tra việc ghi/đọc file JSON trong \texttt{data/profiles}.
    \item \textbf{TC07 – Xây dựng release}: Chạy bản release, xác nhận frontend được tải từ \texttt{frontend/dist}.
\end{itemize}

\subsection{Kết quả kiểm thử thực tế}
\label{subsec:test_results}
\begin{table}[h!]
\centering
\caption{Bảng tổng hợp kết quả kiểm thử}
\label{tab:test_results}
\begin{tabularx}{\textwidth}{|L{4.5cm}|C{3.0cm}|Y|}
\hline
\thcell{Kịch bản} & \thcell{Kết quả} & \thcell{Ghi chú} \\ \hline
TC01 – Khởi động & Thành công & URL đúng (dev) \\ \hline
TC02 – Bắt đầu trò chơi & Thành công & Puzzle nhận được qua WebChannel \\ \hline
TC03 – Đáp án đúng & Thành công & Điểm tăng, thời gian giảm \eqref{eq:score} \\ \hline
TC04 – Đáp án sai & Thành công & Không tăng điểm, thông báo lỗi \\ \hline
TC05 – Hết thời gian & Thành công & Gửi \texttt{gameOver}, lưu hồ sơ JSON \\ \hline
TC06 – Lưu/ tải hồ sơ & Thành công & File JSON hợp lệ, không lỗi parse \\ \hline
TC07 – Release build & Thành công & Frontend tĩnh được hiển thị \\ \hline
\end{tabularx}
\end{table}
Tất cả các kịch bản đều vượt qua, chứng minh tính ổn định của kiến trúc kết nối WebChannel và quá trình biên dịch.

% -------------------------------------------------
% Kết luận triển khai
% -------------------------------------------------
\section{Kết luận triển khai}
\label{sec:implementation_conclusion}
Quá trình triển khai WordSprint kết hợp CMake/Qt cho phần backend và Vite/React cho phần frontend đã đạt được mục tiêu:
\begin{itemize}
    \item Quản lý nguồn và phụ thuộc một cách nhất quán qua CMake.
    \item Nhúng ứng dụng web trong Qt WebEngineView, cho phép giao tiếp hai chiều qua WebChannel.
    \item Kiểm thử toàn diện từ cấp độ unit đến end‑to‑end, đảm bảo chức năng cốt lõi và xử lý ngoại lệ.
    \item Kiến trúc mở rộng: các lớp puzzle mới có thể được thêm bằng mẫu Factory mà không thay đổi giao diện.
\end{itemize}
Với quy trình build và thử nghiệm đã được chuẩn hoá, dự án sẵn sàng cho việc đóng gói và phát hành trên nền Windows.
